\section{Compact supports on independent extension tuples}
\label{sec:independent-tuples}

Fix $m\geq 1$, and fix an embedding in $C=\Cp$ of the unramified
extension $F_m/\Qp$ of degree $m$.  In this section write
\[
 \mathcal N_m=\mathbb G_{a,C}^{\diamond}/\underline{F_m},
 \qquad
 Y_s=(\mathcal N_m^s)_{\mathrm{ind}},
 \qquad Y_0=\Spa(\Cflat).
\]
Here independence means $\Qp$-linear independence on every geometric
fiber.  This is the additive presentation of the negative
Banach--Colmez space in \cref{thm:tower-comparison}; an action on
$\mathcal N_m$ need not lift to a linear action on the displayed
untilted affine coordinate.  We will use the natural commuting actions
of $P_m$ and $M_s=\GL_s(\Zp)$ on $Y_s$.

All cohomology is on the $v$-site.  We use the support convention of
\cref{sec:kummer}: for an increasing, interleaved system of
quasicompact bounds $K\Subset X$,
\begin{equation}
 \label{eq:ind-support-convention}
 \RGamma_c(X,\mathscr E)
 =\colim_K\operatorname{fib}
   \bigl(\RGamma(X,\mathscr E)
       \longrightarrow\RGamma(X\setminus K,\mathscr E)\bigr).
\end{equation}
Bounds may be replaced by slightly smaller closed and slightly larger
open bounds.  In particular, the formula does not require a choice at
a higher-rank boundary point.  The sheaf $\mathscr E$ will be either
$\Fp$ or $\cO^\flat$.

A compact profinite parameter space $\mathcal D$ is retained throughout
each support calculation.  Supports are uniform in this parameter.
Set
\[
 R_{\mathcal D}=C(\mathcal D,\Cflat),\qquad
 E_{\mathcal D}=C(\mathcal D,\Fp).
\]
The second ring consists of locally constant functions.  The first is
equipped with its supremum norm.  When $\mathcal D$ is infinite,
finiteness over $R_{\mathcal D}$ will never mean finite dimension over
$\Cflat$.

\subsection{The translation quotient and its support maps}

We first compute the space before deleting dependent tuples.  The
noncompactness of $F_m$ makes it necessary to keep track of the
support in the quotient.

\begin{lemma}[Lattice transitions]
\label[lemma]{lem:ind-lattice-transitions}
Let $L$ be a lattice of rank $e$ in a finite-dimensional $\Qp$-vector
space, and let $L'=p^{-1}L$.  With trivial characteristic-$p$
coefficients, continuous cohomology is the exterior algebra on
$\Hom_{\cts}(L,\Fp)$.  Under the identifications from a lattice basis,
\begin{enumerate}[label=\textup{(\roman*)}]
 \item restriction from $L'$ to $L$ is zero in positive degrees and
       is the identity in degree zero;
 \item corestriction from $L$ to $L'$ is zero in degrees less than $e$
       and identifies the top-degree orientation lines.
\end{enumerate}
The assertions hold with coefficients $R_{\mathcal D}$ and with
relative base complexes in place of $\Fp$.
\end{lemma}

\begin{proof}
A completed Koszul resolution on the $e$ commuting generators of $L$
computes its continuous cohomology.  With trivial coefficients its
differentials vanish, giving the asserted exterior algebra.  One can
compute the maps first over $\Zp$.  If a basis of $L'$ is $v_1,\ldots,v_e$,
then a basis of $L$ is $pv_1,\ldots,pv_e$.  Restriction in degree $i$
is multiplication by $p^i$ on the corresponding exterior basis.
The identity
$\operatorname{res}\operatorname{cor}=[L':L]=p^e$ and torsion-freeness
then give multiplication by $p^{e-i}$ for corestriction.  Reducing
modulo $p$ proves both assertions.  The same finite Koszul complexes,
tensored with $R_{\mathcal D}$ or with a base complex, compute the
relative assertions.  Thus no infinite product of coefficient spaces
is introduced in this computation.
\end{proof}

\begin{proposition}[The whole extension space]
\label[proposition]{prop:ind-whole}
For every $s\geq 0$ there are natural comparisons
\begin{align}
 \RGamma_c(\mathcal N_m^s\times\underline{\mathcal D},\Fp)
   &\simeq E_{\mathcal D}[-(m+2)s],
       \label{eq:ind-whole-fp}\\
 \RGamma_c(\mathcal N_m^s\times\underline{\mathcal D},\cO^\flat)
   &\simeq R_{\mathcal D}[-(m+2)s].
       \label{eq:ind-whole-flat}
\end{align}
A geometric Kummer orientation is understood.  The comparison from
the first formula to the second is extension of constants.  It is
compatible with continuous compact families of automorphisms and
with all the support and lattice maps used below.
\end{proposition}

\begin{proof}
Consider first $s=1$.  Let $D_a$ be the closed disc of radius
$R_a=R_0|p|^{-a}$ in the untilted affine line, and put
\[
 L_a=\{b\in F_m:|b|\leq R_a\},\qquad
 B_a=\operatorname{image}(D_a\longrightarrow\mathcal N_m).
\]
Choose $R_0$ between adjacent nonzero norm values of $F_m$ and
interleave these discs with slightly larger ones.  Then
$L_{a+1}=p^{-1}L_a$, the $B_a$ exhaust the quotient, and these bounds
are cofinal for \cref{eq:ind-support-convention}.  Indeed every point
has a representative of finite norm, and every quasicompact bound is
contained in one of the increasing slightly-open images.

The inverse image of $B_a$ in the affine line is
\[
 \coprod_{b\in F_m/L_a}(D_a+b).
\]
The displayed pieces are disjoint and locally finite: on any bounded
disc only finitely many cosets occur.  Consequently sections with
this support form the product of the sections with the indicated
disc supports.  With the translation action this is the continuous
coinduction from $L_a$.  This assertion also applies to derived
sections: on bounded charts it is finite disjoint excision, and an
increasing exhaustion gives the displayed product.  Descent through
the translation torsor and Shapiro's map therefore identify the
quotient-support complex with
\begin{equation}
 \label{eq:ind-shapiro-support}
 \RGamma\bigl(L_a,
    \RGamma_{D_a}(\mathbb G_{a,C}^{\diamond},\mathscr E)\bigr).
\end{equation}
Here Shapiro's map can be read directly on the bar complexes: choose
the continuous coset section for the open subgroup $L_a\subset F_m$
and insert the chosen representatives.  The inverse inserts the
same representatives in the coinduced function.  Their homotopies
are the usual insertion homotopies.  Thus this use of Shapiro does
not impose compact inverse image support under the noncompact
translation cover.  After Shapiro all bar parameters belong to
compact powers of $L_a$.

Put
$A_a=\RGamma_{D_a}(\mathbb G_{a,C}^{\diamond},\mathscr E)$.
The canonical maps
\[
 A_a\longrightarrow
 \RGamma_c(\mathbb G_{a,C}^{\diamond},\mathscr E)
       \simeq \mathscr E_{\mathrm{base}}[-2]
\]
are the oriented comparison maps of \cref{prop:kum-line}.
Their cones become cohomologically zero under a sufficiently large
fixed ratio of support radii.  This holds before group cohomology,
with arbitrary compact parameters.  To use it in
\cref{eq:ind-shapiro-support}, first enlarge the geometric disc while
retaining the smaller group $L_a$.  The cone map is zero on the
geometric cohomology rows with parameters $L_a^j$.  Only afterwards
group the translated copies in a larger disc.  Each grouping is a
finite sum, indexed by the appropriate lattice quotient, and hence
preserves the vanishing of the cone map on those rows.  This order
avoids any requirement that a chosen representative of a Kummer
class remain invariant when its center lies near the new boundary.

Under increasing $a$, the old disc supports inside a new disc are
indexed by $L_{a'}/L_a$.  The support map is the sum of their translated
extension-of-support maps.  On the oriented base classes it is
precisely corestriction.  We therefore obtain
\begin{equation}
 \label{eq:ind-quotient-corestriction}
 \RGamma_c(\mathcal N_m,\mathscr E)
 \simeq
 \colim_{a,\,\operatorname{cor}}
       \RGamma(L_a,\mathscr E_{\mathrm{base}})[-2].
\end{equation}
To justify the comparison on total complexes, fix a total degree.
The geometric and positive bar complexes are uniformly bounded
below, so only finitely many bar degrees enter.  Take a common
radius enlargement in those degrees.  The cone has zero colimit on
every relevant geometric row, and the resulting spectral sequence
gives zero colimit in that total degree.  Filtered colimits are exact
on the underlying vector spaces.  Thus this reasoning proves
\cref{eq:ind-quotient-corestriction} without interchanging an
uncontrolled limit and an unbounded totalization.

\Cref{lem:ind-lattice-transitions} leaves only the degree-$m$
orientation line in \cref{eq:ind-quotient-corestriction}.  This proves
the case $s=1$.  For $s$ coordinates use product boxes in the affine
space and the lattice $L_a^s$ of rank $ms$.  Product boxes with a
common increasing bound are cofinal.  The product Kummer calculation
gives degree $2s$, and the same lattice calculation leaves degree
$ms$, proving both displayed formulas.

All these calculations are relative.  On an arbitrary perfectoid
base the phrase ``base class'' means its base complex, not that
arbitrary perfectoid bases have acyclic $\Fp$-cohomology.  For the
external base $\underline{\mathcal D}\times\Spa(\Cflat)$ these base
complexes are $E_{\mathcal D}$ and $R_{\mathcal D}$, respectively,
by \cref{prop:kum-line}.  The Kummer operators are linear over the
retained coefficient ring and their bounds are independent of the
compact parameter.  This proves the parameter assertion and its
compatibility with extension of constants.  The constructions use
canonical restriction and support maps, so their identifications
are natural under compact continuous families of automorphisms.
\end{proof}

\subsection{Persistent lattices and graph neighborhoods}

We now give the uniform neighborhood construction used for rank
deletion.  Its persistence under all further refinements is needed
both for compact supports and, later, for ordinary functions.

\begin{proposition}[Uniform graph tubes]
\label[proposition]{prop:ind-tubes}
Let a quasicompact perfectoid chart over $\Cflat$ carry untilted
coordinates $z_1,\ldots,z_t$ whose classes in $\mathcal N_m$ are
independent.  Restrict to a bounded part of this chart and to a
compact open slope patch in $\Qp^t$.  After a finite refinement of
the base by quasicompact charts, the following constructions are
available on each member of the refinement.
\begin{enumerate}[label=\textup{(\roman*)}]
 \item There are radii $0<\alpha<\beta$ and a fixed lattice
       $L\subset \Qp^t\oplus F_m$ such that, on every geometric fiber,
       \begin{equation}
       \label{eq:ind-persistent-lattice}
       L=\{w:|\ell_z(w)|\leq\alpha\}
        =\{w:|\ell_z(w)|<\beta\},
       \qquad
       \ell_z(a,b)=\sum_i a_i z_i+b.
       \end{equation}
       The same base chart has lattice $p^jL$ for the radii
       $\alpha|p|^j$ and $\beta|p|^j$, for every integer $j$.
 \item Write $L_a=\operatorname{pr}_{\Qp^t}(L)$ and
       $L_b=L\cap F_m$.  At every sufficiently large $j$, the slope
       patch is a finite union of $p^jL_a$-cosets.  Its graph in
       $\mathcal N_m$ has disjoint tubular neighborhoods indexed by
       those cosets.  Each tube is the image of an actual relative
       disc, centered at a section $\sum_i a_i z_i$ for a constant
       representative $a$, modulo the compact translation lattice
       $p^jL_b$.
 \item The closed tubes of radius $\alpha|p|^j$ and the open tubes
       of radius $\beta|p|^j$ are interleaved and cofinal among
       neighborhoods of the graph family.  The construction works
       on a slightly enlarged independent-base bound and on its
       support collars, and it retains an arbitrary compact
       profinite parameter $\mathcal D$.
\end{enumerate}
For finitely many added coordinates and compact open patches of
matrices of slopes, take products of these constructions.  Their
finite base refinements can be chosen simultaneously.
\end{proposition}

\begin{proof}
Let $E=\Qp^t\oplus F_m$, endowed with a fixed $\Qp$-norm.
Independence says that $\ell_z:E\to C'$ is injective at each
geometric point with untilt field $C'$.  Its unit sphere is compact.
On a bounded base chart there is a uniform upper bound
$|\ell_z(w)|\leq M\|w\|$.  At a fixed geometric point there is also
a positive lower bound on the unit sphere.  Choose a sufficiently
fine finite net in that sphere.  The upper bound shows that the
error in replacing a coefficient by a net point is smaller than the
chosen lower bound.  Thus finitely many strict analytic inequalities
on those net points give the same lower bound on an open base
neighborhood.  Quasicompactness gives a finite refinement with
uniform bounds
\begin{equation}
 \label{eq:ind-uniform-norm}
  c\|w\|\leq |\ell_z(w)|\leq M\|w\|,
       \qquad c>0.
\end{equation}
One may refine by rational charts after introducing the strict
margins.  A strict rank-one inequality persists at all the
higher-rank points in such a chart, which is why interleaved radii
are used.

At one geometric point choose a compact coefficient annulus large
enough to contain all coefficients whose images could have norm
near a proposed radius.  This is possible by
\cref{eq:ind-uniform-norm}.  On this annulus the image avoids zero.
The nonarchimedean norm is locally constant away from zero; hence
it takes only finitely many values on the compact annulus.  Choose
$\alpha<\beta$ in a gap between those values.  The preimage of the
disc is an open compact additive subgroup of $E$, and therefore a
lattice $L$.  A finite coefficient net on the annulus, with the
upper bound in \cref{eq:ind-uniform-norm}, makes the inequalities
defining this same $L$ persist on an open base neighborhood.
Outside the annulus the upper and lower bounds already determine
membership.  Taking a finite base refinement proves
\cref{eq:ind-persistent-lattice} on the entire refined chart.
Finally
$\ell_z(p^jw)=p^j\ell_z(w)$, proving persistence for every depth on
the same chart.

Both $L_a$ and $L_b$ are lattices.  In addition, the exact sequence
\[
 0\longrightarrow L_b\longrightarrow L
     \longrightarrow L_a\longrightarrow0
\]
splits as a sequence of finite free $\Zp$-modules.  Fix such a
splitting once on the chart.  Its $F_m$-component gives a continuous
translation choice for every member of a slope cluster.  In
particular, if $a'-a\in p^jL_a$, there is $b\in F_m$ with
$(a'-a,b)\in p^jL$, so that
\[
 \left|\sum_i(a'_i-a_i)z_i+b\right|
       \leq\alpha|p|^j.
\]
Thus the graph of that entire slope coset lies in the tube centered
at $\sum_i a_i z_i$.  The splitting ensures this assertion holds
for the relative family, not only on individual fibers.

Conversely, two such quotient tubes of open radius
$\beta|p|^j$ can intersect only if, for some $b\in F_m$,
$|\ell_z(a-a',b)|<\beta|p|^j$.  By
\cref{eq:ind-persistent-lattice} this forces
$a-a'\in p^jL_a$.  Distinct slope cosets therefore give disjoint
tubes.  The stabilizer of any one covering disc is
$\{b\in F_m:|b|<\beta|p|^j\}=p^jL_b$, with the same group for the
closed radius.  Its full inverse image under the translation cover
is the disjoint union indexed by $F_m/p^jL_b$.

Because the slope patch is compact and open, it is a finite union
of $p^jL_a$-cosets for every sufficiently large $j$.  Choose
representatives at each depth.  If a refined representative lies
over a coarser one, the splitting just chosen supplies a constant
translation making its disc lie in the coarser disc.  Thus all
maps between the tubes are maps between actual relative analytic
discs.  On passing to the quotient they are independent of that
translation choice.
For a child disc, the ultrametric inequality also permits its parent
disc to be centered at the child's section without changing the
parent disc: the two centers differ by at most the parent radius.
Thus the child-to-parent maps are the concentric support maps of
\cref{prop:kum-line}.  Taking a fixed number of further depths makes
their radius ratio at least the constant required there.

We spell out neighborhood cofinality.  A rational neighborhood of
an analytic section contains a relative disc of some positive
radius around that section after a quasicompact base refinement.
Indeed write its finitely many defining inequalities using an
analytic denominator.  On a suitable base chart that denominator
has a positive lower norm bound along the section.  Taylor
estimates on a bounded ambient disc bound all the changes in the
numerators and denominator by a constant times the distance to
the section.  For a sufficiently small radius, the denominator
keeps its norm and the defining inequalities remain true, by the
ultrametric inequality.  This argument also works when an
inequality holds with equality at the section.

Apply this observation to an open neighborhood of the compact
graph family.  The product of a quasicompact base chart and a
compact profinite coefficient patch is quasicompact; its topology
has the usual basis of base opens and clopen parameter patches.
A finite refinement therefore supplies finitely many positive
disc radii.  A sufficiently deep common $p^jL$-refinement, and the
smaller of those finitely many radii, put all its tubes in the
given neighborhood.  Their intersection is exactly the graph:
by \cref{eq:ind-uniform-norm}, a bounded convergent sequence of
centers has a convergent coefficient subsequence, and the
coefficient patch is closed.  This proves both cofinality and the
intersection assertion.

For support fibers, make this construction first on a slightly
larger independent-base bound and on the finitely many rational
pieces of its intervening collars.  These bounds retain the
strict lower bound in \cref{eq:ind-uniform-norm}.  The first $t$
tuple coordinates are unchanged by forming a normal tube, so the
tubes remain in the rank-at-least-$t$ locus.  Bounded matrix slopes
bound all remaining tuple coordinates.  Consequently product
bounds in the base and the compact slope patch are cofinal with
the tuple bounds in this chart.  Conversely, a tuple bound
separated from lower rank has an independent-base bound of this
form, by \cref{eq:ind-uniform-norm}.  This proves the assertion
about collars and supports.

Finally an external compact parameter may be included in the
base chart and in its supremum norm, or treated by finite clopen
refinement.  All finite-net estimates retain the same bounds.
Products use finitely many copies of $E$; taking common refinements
and the maximum of finitely many required radius losses proves
the last assertion.
\end{proof}

\subsection{Rank deletion with neighborhood germs}

For $0\leq t\leq s$ let $X_{s,\geq t}\subset\mathcal N_m^s$ be
the locus of rank at least $t$, and let $Z_{s,t}$ be its rank-$t$
locus.  The latter is closed in $X_{s,\geq t}$: dependence can be
tested on projective coefficient directions, which form a compact
$\Qp$-space.  On a bounded chart the graph inequalities in
\cref{prop:ind-tubes} give the same closedness and show that the
complement is exhausted by bounds with a positive gap from
$Z_{s,t}$.

View a tuple as the linear map $\Qp^s\to\mathcal N_m$ sending
the standard basis to its entries.  Its rank-$t$ coefficient
quotient is $\Qp^s/\ker(f)$; equivalently its coefficient plane
is the annihilator of $\ker(f)$ in the dual space.  We denote
this compact quotient Grassmannian by
$\operatorname{Gr}_t(\Qp^s)$ and use the action induced by
precomposition on tuples.  On its graph charts take the first $t$
independent tuple coordinates as the base.  Each remaining
coordinate has the form
\begin{equation}
 \label{eq:ind-graph-centers}
  \sum_{i=1}^t a_i z_i+b,
       \qquad a_i\in\Qp,\quad b\in F_m,
\end{equation}
after lifting to untilted coordinates.  The normal neighborhoods
are exactly those of \cref{prop:ind-tubes}.

The maps now come from restriction to shrinking neighborhoods.
After descent through the translation cover, they restrict cohomology
to smaller lattice stabilizers. By \cref{lem:ind-lattice-transitions},
positive lattice degrees disappear, while a degree-zero class is
copied into each refined slope cluster. This explains the locally
constant slope families in the germ term below. Earlier, expanding
compact supports used corestriction and retained the top lattice
orientation instead.

\begin{lemma}[The excision term]
\label[lemma]{lem:ind-germ-excision}
There is a natural triangle
\begin{equation}
 \label{eq:ind-rank-triangle}
 \RGamma_c(X_{s,\geq t+1},\mathscr E)
 \longrightarrow \RGamma_c(X_{s,\geq t},\mathscr E)
 \longrightarrow \mathcal G_{s,t}(\mathscr E)
 \longrightarrow
 \RGamma_c(X_{s,\geq t+1},\mathscr E)[1],
\end{equation}
where $\mathcal G_{s,t}$ is computed by shrinking neighborhood
germs of $Z_{s,t}$, with compact support along that locus.
Locally on a compact open Grassmannian graph patch $P$, it is
\begin{equation}
 \label{eq:ind-germ-formula}
  \operatorname{LC}\bigl(P,
           \RGamma_c(Y_t,\mathscr E)\bigr).
\end{equation}
Here $\operatorname{LC}(P,-)$ means the filtered colimit of finite
direct sums over finite clopen partitions of $P$.  With an external
compact parameter $\mathcal D$, the corresponding formula is
\[
 \operatorname{LC}\bigl(P,
   \RGamma_c(Y_t\times\underline{\mathcal D},\mathscr E)\bigr).
\]
At each tube depth only a finite partition of $P$ is used, while
$\mathcal D$ remains in the base complex.  The slope-partition
colimit is algebraic.
\end{lemma}

\begin{proof}
First fix a support bound in $X_{s,\geq t}$ and a slightly larger
bound.  Restrict a section with that support to a neighborhood of
the closed rank-$t$ locus inside the larger bound.  The fiber of
this restriction is a section supported off that neighborhood;
one can replace the neighborhood by a slightly smaller one to
separate all boundary radii.  This is the Mayer--Vietoris support
fiber sequence on the bound and its collar.  As the neighborhood
shrinks, the remaining bounds are cofinal among compact bounds
in $X_{s,\geq t+1}$.  Conversely any such compact bound is
disjoint from a sufficiently small neighborhood of the closed
locus, by the strict separation on a slightly enlarged bound.
Taking the exact filtered colimit of these support triangles and
then increasing the original bound proves
\cref{eq:ind-rank-triangle}.  In particular its third term uses
restriction to shrinking neighborhoods.  It is not an inverse
limit of cohomology supported in shrinking tubes.

To calculate that third term, use \cref{prop:ind-tubes}.  On a
fixed base chart, at depth $j$ there are finitely many slope
clusters and actual relative disc tubes.  Upstairs under the
translation cover, other discs are the disjoint translates of
these discs.  Shapiro reduces their cohomology to the compact
stabilizers $p^jL_b$; this is the same explicit bar construction
as in \cref{eq:ind-shapiro-support}.

The comparison is induced by actual base inclusions.  On a refined
base chart $S$, write $C_S=\RGamma(S,\mathscr E)$, retaining any
external parameter in $S$.  Pullback from $S$ to each covering disc
is translation-equivariant.  After Shapiro it gives the map
\[
 \bigoplus_{\text{slope clusters at depth }j}
       \RGamma(p^jL_b,C_S)
       \longrightarrow \RGamma(\text{union of tubes},\mathscr E).
\]
These maps commute with refinements and with the restriction maps
on the base and its support collars.

Restriction through a sufficiently large fixed ratio of tube
radii has pure base image on geometric cohomology, by
\cref{prop:kum-line}.  The radius ratio is independent of the
center and of compact parameters.  The finite refined cover of
\cref{prop:ind-tubes} works for every depth; hence one further
depth works simultaneously for its charts, its support collars,
and all the compact stabilizer-bar parameters in any fixed total
degree.  It follows that the geometric cone rows vanish before
applying the group maps.  The group maps here are restrictions
to $p^{j'}L_b\subset p^jL_b$.  By
\cref{lem:ind-lattice-transitions}, they kill every positive
lattice-cohomology degree and preserve the base degree zero.
The resulting neighborhood-germ complex is therefore a copy of
the base complex for each slope cluster.  When a cluster splits,
restriction replicates its class in each new cluster.

This argument is relative over the independent base, rather
than just a calculation on its geometric fibers.  More
explicitly, the centers in \cref{eq:ind-graph-centers} are
sections on entire refined charts, so the Kummer homotopies are
linear over their base rings.  Pullback of a base class realizes
the comparison on each tube.  On overlaps the comparisons are
the restrictions of the same class; a further interleaving of
tubes compares different choices.  For descent over a support
bound and its collar, choose a $v$-hypercover by quasicompact
perfectoid charts in every degree needed.  Fiber products are
quasicompact after the bounded translation reduction.  In a
fixed total degree only finitely many of these degrees occur.
One common further tube depth kills the cone rows there, so the
bounded-below descent spectral sequence proves the relative
comparison.  Applying the support fiber on the base gives
\cref{eq:ind-germ-formula}.

There are only finite sums at each partition.  Their exact
filtered colimit is $\operatorname{LC}(P,-)$, and an external
compact parameter enters the coefficient ring of these finite
sums.  Thus the construction allows continuous coefficients in
$\mathcal D$ but does not allow a single germ to use
unboundedly fine slope partitions with no common neighborhood.
The comparisons agree under restriction to the actual graphs,
which also proves their naturality under changes of graph
chart and changes of lifts.
\end{proof}

\begin{theorem}[Scalar comparison and admissibility]
\label{thm:ind-comparison}
For every $s\geq0$ and every compact profinite $\mathcal D$,
extension of constants induces natural quasi-isomorphisms
\begin{align}
 E_{\mathcal D}\otimes_{\Fp}\RGamma_c(Y_s,\Fp)
   &\xrightarrow{\ \sim\ }
       \RGamma_c(Y_s\times\underline{\mathcal D},\Fp),
       \label{eq:ind-fp-parameters}\\
 R_{\mathcal D}\otimes_{\Fp}\RGamma_c(Y_s,\Fp)
   &\xrightarrow{\ \sim\ }
       \RGamma_c(Y_s\times\underline{\mathcal D},\cO^\flat).
       \label{eq:ind-algebraic-comparison}
\end{align}
The tensors are algebraic.  The comparisons retain uniform
support and all compact bar parameters.

Each $H_c^i(Y_s,\Fp)$ is a countable smooth admissible
$M_s$-module.  It has a smooth commuting $P_m$-action.  Smooth
means that each vector has open stabilizer; admissible means
that every open subgroup of $M_s$ has finite-dimensional fixed
space.
\end{theorem}

\begin{proof}
Induct on $s$.  The case $s=0$ is the coefficient calculation on
the external parameter space.  For the induction step start
with the whole space $X_{s,\geq0}=\mathcal N_m^s$, whose
comparison is \cref{prop:ind-whole}.  Delete the ranks
$t=0,\ldots,s-1$ successively using
\cref{eq:ind-rank-triangle}.  For each such $t$, the germ term
is locally \cref{eq:ind-germ-formula}.  By the inductive
hypothesis it has the scalar comparison in question.  The
functor $\operatorname{LC}(P,-)$ is an exact filtered colimit
of finite direct sums and commutes with algebraic scalar
extension.  For its gluing, use integral frames of the coefficient
quotients: the image of $\Zp^s$ in a $t$-dimensional coefficient
quotient is a free lattice, and its kernel in $\Zp^s$ is a
saturated direct summand.  Dually, the coefficient plane meets
the dual integral lattice in a saturated rank-$t$ summand.
Consequently $M_s$ acts transitively on the Grassmannian, and
integral graph charts give continuous local sections of
$M_s\to M_s/\Pi_t$, where $\Pi_t$ stabilizes the standard plane.
The resulting transition functions on the independent base
lie in $\GL_t(\Zp)$.  Their action on any finite set of
cohomology classes is locally constant by inductive smoothness;
compactness gives a common finite clopen refinement.  A finite
disjoint clopen trivializing cover computes global sections by
finite sums of the exact functors $\operatorname{LC}(P,-)$.
Thus the germ comparison glues with its algebraic finite-span
structure on the compact Grassmannian.  Exactness of scalar extension
over $\Fp$ and the support triangles prove
\cref{eq:ind-fp-parameters,eq:ind-algebraic-comparison} for the
final complement $Y_s$.

The proof took place on uniform compact parameters before the
support colimit and before each bounded-below totalization,
as established in \cref{lem:ind-germ-excision}.  It therefore
gives the asserted comparisons on actual coefficient
complexes.  In particular it does not replace continuous
families of unrestricted functions by an algebraic tensor
after taking cohomology.

Smoothness, including the commuting $P_m$-action, now follows
from the parameter comparison.  For a compact acting group
$H$ and a class $v$, pullback by its action gives a family
of classes with parameter $\mathcal D=H$.  Supports remain
uniform: a compact group translates a quasicompact bound
into a quasicompact image contained in a larger bound.
By \cref{eq:ind-fp-parameters}, this family lies in
$C(H,\Fp)\otimes_{\Fp}H_c^i(Y_s,\Fp)$.  It is therefore
a locally constant finite-valued orbit map.  Thus the
stabilizer of $v$ is open.  This proves smoothness for both
commuting actions before we use the category of smooth
representations below.

We next identify the representation in a germ term.  Let
$\Pi_t\subset M_s$ be the stabilizer of the standard
coefficient $t$-plane.  Its action on the independent base is
through its $\GL_t(\Zp)$ Levi factor; the other Levi factor
and the unipotent radical act trivially on the normal
degree-zero germ.  Hence
\begin{equation}
 \label{eq:ind-parabolic-module}
 H^i\mathcal G_{s,t}(\Fp)
   \cong
 \Ind_{\Pi_t}^{M_s} H_c^i(Y_t,\Fp),
\end{equation}
where induction is smooth induction and no support condition
on $M_s/\Pi_t$ is needed because that quotient is compact.
To see the identification, use the integral trivializations
just chosen on the graph charts.  Sections are locally constant choices
of an independent-base class, and transition functions are
the changes of basis of that plane.  These are exactly the
equivariant locally constant functions defining the right
side.  The comparison is by restriction to the plane, so
it is compatible with the group action and not just an
isomorphism of vector spaces.

Smooth parabolic induction in \cref{eq:ind-parabolic-module}
preserves admissibility.  Indeed an open subgroup $H\subset
M_s$ has finitely many orbits on $M_s/\Pi_t$.  For each orbit,
the stabilizer $H\cap g\Pi_tg^{-1}$ maps onto an open
subgroup of the $\GL_t(\Zp)$ Levi factor.  An $H$-fixed
induced function is determined by a vector fixed under that
open subgroup.  There are finitely many such choices, each
finite-dimensional by induction.

For completeness, admissible smooth representations of a
compact $p$-adic analytic group over $\Fp$ are closed under
subobjects, quotients, and extensions.  Take continuous
duals into $\Fp$.  A smooth module is admissible precisely
when its compact dual is finitely generated over the
completed group algebra: finite coinvariants for an open
uniform pro-$p$ subgroup give finite generation by compact
Nakayama, and the converse follows by taking coinvariants
at open subgroups.  The completed group algebra is
Noetherian, since the uniform subgroup algebra has its
polynomial associated graded ring and the full algebra is
finite over it.  Finite generation is therefore closed
under submodules, quotients, and extensions.  Dualizing
proves the assertion for smooth modules.

The whole-space cohomology in \cref{prop:ind-whole} is finite.
The long exact sequences of
\cref{eq:ind-rank-triangle}, \cref{eq:ind-parabolic-module},
and the preceding closure property now prove admissibility
for $Y_s$.  Countability follows in the same induction:
the compact $p$-adic Grassmannian has a countable clopen
basis, so locally constant functions on it with values in
a countable $\Fp$-space form a countable space, and the
support exact sequences preserve countability.

\end{proof}

\subsection{Finite equivariant compact-support cohomology}

\begin{lemma}
\label[lemma]{lem:ind-admissible-cohomology}
Let $H$ be a compact $p$-adic analytic group and let $V$ be a
countable smooth admissible $\Fp[H]$-module.  Then
$H^a_{\cts}(H,V)$ is finite-dimensional for every $a$.
Let $H$ act trivially on the external parameter
$\mathcal D$ and on $R_{\mathcal D}$, and diagonally through
$V$ on $R_{\mathcal D}\otimes V$.  For the coefficient
topology supplied by uniform finite-span families, there is a
natural isomorphism
\[
 H^a_{\cts}(H,R_{\mathcal D}\otimes_{\Fp}V)
 \cong R_{\mathcal D}\otimes_{\Fp}H^a_{\cts}(H,V).
\]
\end{lemma}

\begin{proof}
Let $V^\vee=\Hom(V,\Fp)$ be the compact dual.  It is
finitely generated over $\Fp[[H]]$, as in the proof of
\cref{thm:ind-comparison}.  Noetherianity gives a resolution
of $V^\vee$ with finitely generated free terms, finite in
each degree.  Tensoring it with the trivial module $\Fp$
shows that $\Tor_a^{\Fp[[H]]}(\Fp,V^\vee)$ is finite in
each degree.  Continuous duality identifies this with the
dual of $H^a_{\cts}(H,V)$.

For scalar extension use the finite free resolution of the
trivial module supplied by \cref{lem:cts-comparison}.
Exhaust $V$ by finite-dimensional $H$-stable subspaces.
Such an exhaustion exists because
$V$ is countable and every smooth vector has finite orbit
under the compact group.  Give
$R_{\mathcal D}\otimes V$ the union of the complete
finite-span steps so obtained.  Its actual bounded-step
bar cochains are
\[
 C^a_{\cts}(H,R_{\mathcal D}\otimes V)
   =R_{\mathcal D\times H^a}\otimes_{\Fp}V.
\]
Indeed this is true on every finite span and then on their
union.  These are exactly the row families given by
\cref{thm:ind-comparison} with parameter
$\mathcal D\times H^a$, compatibly with every bar face.
The finite-type completed
resolution and its comparison with continuous cochains
apply by \cref{lem:cts-comparison}: each vector belongs to
a finite continuous representation, and compact families
are retained in one such step.  In every degree the
resolution calculation is a finite power of this coefficient
module.  A matrix entry in $\Fp[[H]]$ acts on a finite
$H$-stable span through a finite quotient of $H$, so its
action after extension is $R_{\mathcal D}$-linear scalar
extension of its action on $V$.  The differential is therefore
scalar extension of the differential over $\Fp$.
Exactness of tensoring over $\Fp$ gives the displayed
isomorphism.  These are the same finite-span families as
in \cref{thm:ind-comparison}, not a newly imposed completion
or topology on its cohomology.
\end{proof}

\begin{corollary}[Finite equivariant cohomology]
\label[corollary]{cor:ind-finite}
Let $1\leq m<n$, $r=n-m$, and let $U\subset
P_m\times\GL_r(\Zp)$ be a sufficiently deep product open
subgroup as in \cref{thm:tower-comparison}.  Let $Z_U$ be
the perfected analytic cover associated with $B_U$ over
\[
 0<|x|<1,\qquad |y_i|<1.
\]
Then
\[
 H^j\RGamma\bigl(G,\RGamma_c(Z_U,\cO^\flat)\bigr)
\]
is finite-dimensional over $\Cflat$ for every integer $j$.
With an external compact profinite parameter $\mathcal D$,
the corresponding groups are obtained by extension to
$R_{\mathcal D}$ and are finite over that ring.
\end{corollary}

\begin{proof}
Write $U=U_m\times U_r$.  The two compact frame torsors
and \cref{lem:tower-supports} identify the complex in the
assertion with
\[
 \RGamma\bigl(U,\RGamma_c(Y_r,\cO^\flat)\bigr).
\]
Indeed, on the completed full marked tower the actions
commute.  Descent first by $U$ and then by $G$, or in the
opposite order, gives the two displayed descriptions.
Compact torsors preserve the cofinal systems of support
bounds.  The uniform-parameter comparisons of
\cref{thm:ind-comparison} apply to their bar diagrams.
Bounded-below totalization in each degree therefore
justifies this interchange of descents with the support
exhaustion.

By \cref{thm:ind-comparison} the geometric cohomology rows
are scalar extensions of the smooth admissible modules
$H_c^b(Y_r,\Fp)$ for $M_r$.  Their $U_r$-cohomology is
finite in each degree by
\cref{lem:ind-admissible-cohomology}.  The residual
$U_m$-action on these finite $\Fp$-spaces is smooth, since
the original commuting action is smooth and the
finite-type cochain comparison retains its compact
parameters.  A finite-type resolution for $U_m$ then has
finite-dimensional terms with these finite coefficients;
its cohomology is finite in each degree as well.

The successive first-quadrant spectral sequences have only
finitely many terms in any prescribed total degree.
Their entries, and hence their abutments, are finite over
$\Fp$ before scalar extension.  The same finite-type
comparisons and \cref{eq:ind-algebraic-comparison} identify
the $\cO^\flat$ answers with extension to $\Cflat$, or to
$R_{\mathcal D}$ when the parameter is retained.  This
proves the assertion.  No finite cohomological dimension
of a torsion-containing full stabilizer is used here.
\end{proof}

\subsection{Ordinary functions and inverse support limits}

The following ordinary degree-zero assertion is different
from the neighborhood-germ comparison.  It will identify
an invariant unit occurring in the Cartier normalization.

\begin{corollary}[Functions on the independent locus]
\label[corollary]{cor:ind-functions}
For every $s\geq0$,
\begin{equation}
 \label{eq:ind-ordinary-functions}
 \Gamma(Y_s,\cO^\flat)=\Cflat,
 \qquad
 \Gamma(Y_s\times\underline{\mathcal D},\cO^\flat)
       =R_{\mathcal D}.
\end{equation}
The identifications are by pullback of scalars.
For $s\geq1$, projection $Y_s\to Y_{s-1}$ has relative
degree-zero direct image $\cO^\flat$.

Consequently a $G$-invariant algebraic function in a fixed
finite marked coefficient model becomes a scalar after
pullback to the completed full marking tower.  If the
model and its chosen component are defined over the
finite field $\kk$, that scalar belongs to $\kk$.
In particular this applies to invariant algebraic units.
\end{corollary}

\begin{proof}
For $s\geq1$, work locally over an affinoid perfectoid
chart $S$ of $Y_{s-1}$ carrying untilted lifts
$z_1,\ldots,z_{s-1}$.  Such charts exist by the
translation presentation.  The inverse image of the last
independent coordinate in the untilted affine line is
the complement of the closed graph family
\[
 A_z=\left\{\sum_{i=1}^{s-1}a_i z_i+b:
           a_i\in\Qp,\ b\in F_m\right\}.
\]
Refine the bounded independent-base chart using
\cref{prop:ind-tubes}.  Choose $\rho$ strictly between its
persistent radii $\alpha<\beta$ and exhaust the output
coordinate by discs
\[
 X_a=\{|T|\leq\rho|p|^{-a}\},\qquad a\geq0.
\]
The centers in $X_a$ correspond exactly to the compact open
coefficient lattice $p^{-a}L$, and all have norm at most
$\alpha|p|^{-a}$.  They therefore have a separated outer
collar inside $X_a$.  This uses the persistence for negative
as well as positive powers of $p$; the same linear identity
proves both.  At fixed $a$, and sufficiently large $j$, the
closed tubes $T_j$ before translation quotient are the finite
disjoint discs indexed by $p^{-a}L/p^jL$, with radius
$\alpha|p|^j$ and centers $\ell_z(w)$ at constant coefficient
representatives $w$.  Their open companions have radius
$\beta|p|^j$ and lie inside the same collar.  By
\cref{prop:ind-tubes} these companions are a cofinal
neighborhood system, and their complements admissibly exhaust
$X_a\setminus A_z$.  On the finite base refinement the
exhaustions just chosen have common enlargements, so they
compute the same support fibers on overlaps.

Let $X$ denote this bounded ambient chart with its
chosen collar, and put
\[
 C_j=\operatorname{fib}\bigl(
       \RGamma(X,\cO^\flat)
        \longrightarrow
       \RGamma(X\setminus T_j,\cO^\flat)\bigr).
\]
The maps are $C_{j'}\to C_j$ for $j'\geq j$:
\emph{smaller support maps to larger support}.  A small
cluster is included in its containing larger disc;
the transition sums these maps when several clusters
are included in one larger cluster.  By
\cref{lem:kum-low-support}, after a sufficiently large
fixed ratio of these radii the maps on $H^0$ and
$H^1$ are zero.  This is vanishing of the entire
cohomology-group map, not elementwise eventual
vanishing.  The uniform persistent lattice charts
and support collars are precisely the hypotheses
needed for that lemma.  Finite sums of the cluster
maps retain the vanishing.  On a finite descent
cover one may compose finitely many further
refinements to kill the finite filtration in these
two total degrees.  Hence both low-degree inverse
systems are pro-zero on the whole bounded chart.

Since complements exhaust $X\setminus A_z$, limits
preserve the support fiber and give
\begin{equation}
 \label{eq:ind-closed-support-limit}
 \RGamma_{A_z}(X,\cO^\flat)
       \simeq \Rlim_j C_j.
\end{equation}
The complexes $C_j$ have no negative cohomology.
The Milnor exact sequence for this countable tower
therefore shows
\[
 H^0\RGamma_{A_z}(X,\cO^\flat)
   =H^1\RGamma_{A_z}(X,\cO^\flat)=0:
\]
both ordinary limits of $H^0,H^1$ vanish, and the
possible $\varprojlim^1 H^0$ vanishes because that
system is pro-zero.  No condition on the surviving
degree-two support system is required.

Now increase the output-coordinate bound.  At every fixed bound
the coefficient set $p^{-a}L$ is compact, so the preceding argument
applies.  Use the interleaved open bounds and set
\[
 C_a=\RGamma_{A_z\cap X_a}(X_a,\cO^\flat).
\]
Although the tube refinements may depend on $a$, these support
complexes restrict canonically.  The bounds and their complements
admissibly exhaust the relative affine line and its complement of
$A_z$.  Thus the exhaustion assertion of \cref{lem:kum-low-support}
gives
\[
 \RGamma_{A_z}(\mathbb A_S,\cO^\flat)
       \simeq\Rlim_a C_a.
\]
Each $C_a$ is connective with $H^0(C_a)=H^1(C_a)=0$, so the Milnor
sequence proves the same low-degree vanishing on the whole line.
The two limits have been taken in order: first the tubes at a fixed
output bound, then the canonical support complexes as that bound
increases.

The support triangle and \cref{prop:kum-line} now
give an isomorphism from base functions to the
ordinary functions on the affine line minus
$A_z$.  This is a relative statement on the
perfectoid base charts: all maps and estimates
are linear over the base and agree on overlaps.
Degree-zero descent by the translation group
$F_m$ takes invariants of these base functions,
on which the group acts trivially.  Hence
projection $Y_s\to Y_{s-1}$ has the claimed
relative degree-zero direct image.  Induction,
starting with $Y_0$, proves
\cref{eq:ind-ordinary-functions}.  Retaining
$\mathcal D$ in the same relative calculation
gives its parameter version.

For the last assertion, the finite model injects
into its perfected analytic realization and
then into the full marking tower: analytic
evaluation on the exhausting bounds is
injective, and the marking extension is
faithful.  A $G$-invariant function descends
through the actual torsor of
\cref{thm:tower-comparison}.  By
\cref{eq:ind-ordinary-functions}, its image is
a scalar $\lambda\in\Cflat$.  Apply
coefficient-only $|\kk|$-Frobenius to this
equality back on the tower, whose finite
model is defined over $\kk$.  It fixes the
original algebraic function, whereas it
sends $\lambda$ to $\lambda^{|\kk|}$.
Injectivity gives
$\lambda^{|\kk|}=\lambda$, so
$\lambda\in\kk$.  This argument does not
require the chosen untilted quotient
presentation itself to commute with
coefficient Frobenius.  If a preliminary
finite constant extension was needed to
define the model or its idempotent, the
conclusion is in that finite field.
\end{proof}

\begin{remark}[Three directions of passage]
\label[remark]{rem:ind-limit-directions}
The preceding arguments use three different systems.
Expanding compact support on $\mathcal N_m$ gives a
filtered colimit with \emph{corestriction} between
expanding translation lattices; its top lattice
degree survives.  Neighborhood germs of a deleted
rank stratum give a filtered colimit with
\emph{restriction} to shrinking lattices; only
lattice degree zero survives and slope classes
are replicated.  Cohomology with support in the
actual closed graph is instead
\cref{eq:ind-closed-support-limit}, an inverse
limit of support fibers; the maps sum classes
from smaller supported tubes into larger ones.
Its pro-zero degrees zero and one are what
prove \cref{cor:ind-functions}.  Neither that
ordinary-function conclusion nor
\cref{cor:ind-finite} asserts a comparison
between incomplete algebraic ordinary
cohomology and completed analytic ordinary
cohomology.
\end{remark}

\begin{figure}[tbp]
\centering
\begin{tikzpicture}[
  x=.98\linewidth,y=1cm,
  every node/.style={font=\small,inner sep=1pt},
  heading/.style={anchor=north west,align=left,text width=.28\linewidth},
  outcome/.style={anchor=north west,align=left,text width=.26\linewidth},
  transition/.style={-{Stealth[length=1.8mm]},semithick}
]
\node[heading] at (0,0)
  {\textbf{Expanding supports}\\[3pt]
   $K_a\subset K_{a+1}$};
\node (a0) at (.35,-.42) {$A_a$};
\node (a1) at (.62,-.42) {$A_{a+1}$};
\draw[transition] (a0) -- node[above,font=\footnotesize]
  {corestriction} (a1);
\node[outcome] at (.74,0)
  {$\colim_a A_a\simeq R[-(m+2)]$\\[3pt]
   Lattice degree $m$ survives.};
\draw[black!20] (0,-1.28) -- (1,-1.28);

\node[heading] at (0,-1.65)
  {\textbf{Shrinking neighborhoods}\\[3pt]
   $\mathcal U_{j+1}\subset\mathcal U_j$};
\node (g0) at (.35,-2.07) {$\mathcal G_j$};
\node (g1) at (.62,-2.07) {$\mathcal G_{j+1}$};
\draw[transition] (g0) -- node[above,font=\footnotesize]
  {restriction} (g1);
\node[outcome] at (.74,-1.65)
  {$\colim_j\mathcal G_j\simeq\operatorname{LC}(P,\mathcal B)$\\[3pt]
   Normal degree $0$ survives; slope values replicate.};
\draw[black!20] (0,-3.13) -- (1,-3.13);

\node[heading] at (0,-3.5)
  {\textbf{Shrinking supports}\\[3pt]
   $T_{j+1}\subset T_j$};
\node (e0) at (.35,-3.92) {$E_j$};
\node (e1) at (.62,-3.92) {$E_{j+1}$};
\draw[transition] (e1) -- node[above,font=\footnotesize]
  {support inclusion} (e0);
\node[outcome] at (.74,-3.5)
  {$\Rlim_j E_j\simeq\RGamma_Z(X,\mathscr E)$\\[3pt]
   $H^0=H^1=0$: the low-degree systems are pro-zero.};
\node[align=center,font=\footnotesize,text width=.40\linewidth]
  at (.485,-4.65) {Child-cluster support classes are summed.};
\end{tikzpicture}
\caption{The three limit directions of
\cref{rem:ind-limit-directions}.  For retained coefficients $R$,
$A_a=\RGamma(L_a,R)[-2]$ is the lattice comparison model, with
$L_{a+1}=p^{-1}L_a$.  On a slope patch $P$, the complex
$\mathcal G_j$ uses neighborhood sections with compact support
along the stratum; its base complex
$\mathcal B=\RGamma_c(Y_t,\mathscr E)$ is retained.  The complex
$E_j$ is instead the support fiber on a fixed ambient bound $X$,
with $Z=\bigcap_jT_j$.}
\label{fig:support-directions}
\end{figure}

